% Rubbish 写一篇模板（LaTeX 版）· 2026-11 修订
% 用法：xelatex 编译两次；替换「」内容后，删除本注释与下一段。
% 本模板对应《Rubbish 写作与发布规范》；与 Markdown、Word 模板等价，任选其一。
\documentclass[12pt,a4paper]{article}
\usepackage[margin=3cm]{geometry}
\usepackage{ctex}
\usepackage{graphicx}
\usepackage{booktabs}
\usepackage{amsmath}
\usepackage{hyperref}
\usepackage{caption}
\usepackage{setspace}

\captionsetup{font=small,labelsep=quad}
\onehalfspacing

\title{\textbf{「题名」：一句具体的话，必要时再用副题名交代范围}\\
{\large 研究长文 / 荒诞写作 / 自由创作 —— 三选一，删去其余}}
\author{「笔名 / 稳定署名」\\
{\small 「联系方式（可选）\quad 合作者与分工：查资料 / 做图 / 讨论 / 写作」}}
\date{「日期」}

\begin{document}
\maketitle

\begin{abstract}
「两三百字。写清问题、做法、发现与边界，不写“本文具有重要意义”这类空话。
自由创作可只写故事梗概。」\\[0.8em]
\textbf{话题：}「可选，四到六个；只用于浏览分组，不构成评价」
\end{abstract}

\section{问题}
「现象是什么，为什么值得写。用事实、材料或一个具体的场景开头。」

\section{材料与做法}
「研究长文：数据、样本、文本、访谈、实验或论证框架，让别人能沿着你的路再走一遍。
荒诞写作：设定、材料、编码与推理规则。涉及人的调查，说明知情同意与处理方式。
自由创作：本节可删。」

\begin{equation}
y = a + bx
\label{eq:demo}
\end{equation}

\section{展开}
「分节推进，每节只回答一个问题。图表跟着论点走，不用来撑版面。」

\begin{table}[htbp]
\centering
\caption{「表题（放表上方；三线表，不用竖线）」}
\label{tab:demo}
\begin{tabular}{lcc}
\toprule
「条件」 & 「指标 A」 & 「指标 B」 \\
\midrule
「对照组」 & 0.00 & 0.00 \\
「实验组」 & 0.00 & 0.00 \\
\bottomrule
\end{tabular}
\end{table}

\begin{figure}[htbp]
\centering
% \includegraphics[width=0.8\textwidth]{figure.png}
\caption{「图题（放图下方），末尾写数据来源」}
\label{fig:demo}
\end{figure}

\section{讨论}
「结论能推到哪，推不到哪；有哪些竞争解释，哪些地方还没做完。
荒诞写作请诚实说出“荒诞在哪里”。」

\section{结语}
「四到六句：重申核心发现、点明含义、指出后续方向。不引入新材料。」

\section*{引用与出处}
「列出你真正用到的材料：数据、代码、访谈提纲、图片、档案、现场记录。能公开的给链接；不能公开的，说明限制与获取方式。」

\begin{thebibliography}{99}
\bibitem{ref1} 「作者. 题名. 出处 / 站点 / 书目信息, 年. 链接，或 ISBN、DOI（如有）.」
\bibitem{ref2} 「作者. 书名. 版本. 出版地: 出版社, 年. ISBN.」
\end{thebibliography}

\section*{材料与边界说明}
\textbf{虚构与真人：}「如为虚构、讽刺或思想实验，写明“虚构创作”，并确认未指向可识别个人的隐私与名誉。」

\textbf{AI 辅助：}「写明工具与用途；未使用可写“无”。」

\textbf{数据与代码：}「已公开（给仓库或链接）/ 可索取（向谁要、多久回应）/ 受限（说明原因）。」

\textbf{利益冲突：}「有就写清；没有写“无”。」

\section*{公开范围与许可}
\textbf{公开范围：}「公开（站内与访客可读）/ 仅站内登录可读 / 受限（说明谁能看到）。」

\textbf{使用许可：}「你保留著作权；本站用于展示、索引与备份。转载、翻译、改编请注明来源。」

\textbf{转载与引用：}「他人作品的图片、文字、数据，已注明来源与许可；未获授权的部分不在此列。」

\end{document}
